\documentclass[11pt,a4paper]{book}

% ============================================
% PREAMBLE
% ============================================
\usepackage[utf8]{inputenc}
\usepackage[T1]{fontenc}
\usepackage{lmodern}
\usepackage{amsmath,amssymb,amsthm}
\usepackage{geometry}
\usepackage{titlesec}
\usepackage{fancyhdr}
\usepackage{microtype}
\usepackage{enumitem}
\usepackage{array}

\geometry{
  paperwidth=6.5in,
  paperheight=9.5in,
  textwidth=5in,
  textheight=7.5in,
  top=0.75in,
  bottom=0.75in,
  inner=0.75in,
  outer=0.75in
}

% Theorem environments
\theoremstyle{plain}
\newtheorem{theorem}{Theorem}
\newtheorem{lemma}[theorem]{Lemma}
\newtheorem{corollary}[theorem]{Corollary}
\newtheorem{proposition}[theorem]{Proposition}

\theoremstyle{definition}
\newtheorem{definition}[theorem]{Definition}

\theoremstyle{remark}
\newtheorem{remark}{Remark}
\newtheorem*{note}{Note}

% Custom commands for Cayley's notation
\newcommand{\dd}{\,\mathrm{d}}
\newcommand{\td}{\tilde}
\newcommand{\vp}{\varphi}

% Header/Footer
\pagestyle{fancy}
\fancyhf{}
\fancyhead[C]{\small\nouppercase{\leftmark}}
\fancyfoot[C]{\thepage}
\renewcommand{\headrulewidth}{0.4pt}

% Title formatting
\titleformat{\chapter}[display]
  {\normalfont\Large\bfseries\centering}
  {\chaptertitlename\ \thechapter}{20pt}{\Large}

\titleformat{\section}
  {\normalfont\large\bfseries}
  {\thesection}{1em}{}

\titleformat{\subsection}
  {\normalfont\normalsize\bfseries}
  {\thesubsection}{1em}{}

% ============================================
% DOCUMENT
% ============================================
\begin{document}

% --------------------------------------------
% FRONT MATTER: Pages 1--10
% --------------------------------------------
\frontmatter

% Page 7 -- Half-title
\thispagestyle{empty}
\begin{center}
\vspace*{2in}
{\Large\scshape Mathematical Papers}
\end{center}
\cleardoublepage

% Page 9 -- Title Page
\thispagestyle{empty}
\begin{center}
\vspace*{1.5in}
{\Huge\scshape The Collected\\[0.3em] Mathematical Papers}\\[1.5em]
{\Large\scshape of}\\[1.5em]
{\LARGE Arthur Cayley, Sc.D., F.R.S.,}\\[0.5em]
{\small Sadlerian Professor of Pure Mathematics in the University of Cambridge.}\\[2em]
{\Large Vol.\ I.}\\[2em]
\vfill
{\large Cambridge:}\\
{\normalsize At the University Press.}\\[0.5em]
{\small 1889}\\[0.5em]
{\scriptsize [All Rights reserved.]}
\end{center}
\cleardoublepage

% Page 10 -- Colophon
\thispagestyle{empty}
\begin{center}
\vspace*{3in}
{\small Cambridge:}\\
{\small Printed by C.\ J.\ Clay, M.A.\ and Sons,}\\
{\small At the University Press.}
\end{center}
\cleardoublepage

% --------------------------------------------
% PREFACE: Pages 11--12
% --------------------------------------------
\chapter*{Preface}
\addcontentsline{toc}{chapter}{Preface}
\markboth{Preface}{Preface}

Rather more than a year ago I was requested by the Syndics of the University Press to allow my mathematical papers to be reprinted in a collected form: I had great pleasure in acceding to a request so complimentary to myself, and I willingly undertook the work of superintending the impression of them, and of adding such notes and references as might appear to me desirable.

The present volume contains one hundred papers, originally published in the years 1841--1853 inclusive. These papers are arranged chronologically according to the dates of publication; the first paper is dated February 1841, and the last paper December 1853. The volume includes, in particular, the series of papers on Invariants, the series on Linear Transformations, and the series on Elliptic Functions; as well as various papers on Geometry, on the Theory of Curves, on Definite Integrals, and on other subjects.

\cleardoublepage

% --------------------------------------------
% CONTENTS: Pages 13--19
% --------------------------------------------
\chapter*{Contents}
\addcontentsline{toc}{chapter}{Contents}
\markboth{Contents}{Contents}

{\small [An Asterisk denotes that the Paper is not printed in full.]}

\vspace{1em}
\noindent
\begin{tabular}{@{}r@{\quad}p{3.2in}@{\quad}r@{}}
1. & On a Theorem in the Geometry of Position & \dotfill 1 \\
  & \small Camb.\ Math.\ Jour.\ t.~ii.\ (1841), pp.~267--271 & \\
2. & On the Properties of a Certain Symbolical Expression & \dotfill 7 \\
  & \small Camb.\ Math.\ Jour.\ t.~iii.\ (1841), pp.~62--71 & \\
3. & On Certain Definite Integrals & \dotfill 13 \\
  & \small Camb.\ Math.\ Jour.\ t.~iii.\ (1843), pp.~138--144 & \\
4. & On Certain Expansions, in series of Multiple Sines and Cosines & \dotfill 19 \\
  & \small Camb.\ Math.\ Jour.\ t.~iii.\ (1842), pp.~162--167 & \\
5. & On the Intersection of Curves & \dotfill 25 \\
  & \small Camb.\ Math.\ Jour.\ t.~iii.\ (1843), pp.~211--213 & \\
6. & On the Motion of Rotation of a Solid Body & \dotfill 28 \\
  & \small Camb.\ Math.\ Jour.\ t.~iii.\ (1843), pp.~224--232 & \\
7. & On the Attraction of Ellipsoids & \dotfill 37 \\
  & \small Camb.\ Math.\ Jour.\ t.~iii.\ (1843), pp.~280--284 & \\
8. & On a particular case of the Rectilinear Motion of a System of Bodies & \dotfill 41 \\
  & \small Camb.\ Math.\ Jour.\ t.~iii.\ (1843), pp.~284--287 & \\
9. & On the Intersection of Two Surfaces & \dotfill 44 \\
  & \small Camb.\ Math.\ Jour.\ t.~iii.\ (1843), pp.~290--291 & \\
10. & On the Theory of Analytical Curves & \dotfill 46 \\
  & \small Camb.\ Math.\ Jour.\ t.~iii.\ (1843), pp.~292--294 & \\
11. & Note on the Integration of certain Differential Equations & \dotfill 49 \\
  & \small Camb.\ Math.\ Jour.\ t.~iii.\ (1843), pp.~295--296 & \\
12. & On the Theory of Linear Transformations & \dotfill 51 \\
  & \small Camb.\ Math.\ Jour.\ t.~iii.\ (1843), pp.~302--304 & \\
\end{tabular}

\vfill
\begin{center}
\textit{(Continued in subsequent pages)}
\end{center}
\cleardoublepage

% --------------------------------------------
% CLASSIFICATION: Pages 21--22
% --------------------------------------------
\chapter*{Classification}
\addcontentsline{toc}{chapter}{Classification}
\markboth{Classification}{Classification}

\section*{Geometry}
\begin{itemize}[leftmargin=*,noitemsep]
\item Intersection of curves, 5
\item Motion of solid body, 6, 36, 37, 68, 78
\item Theory of algebraical curves, 10
\item Analytical geometry of $n$ dimensions, 11
\item Curves and developables, 30, 83
\item Geometrical involution, 40
\item Geometry of position, 50, 55, 70, 95, 98
\item Theory of Curves, 53
\item Reciprocity, 61
\item Transmutations of curves, 80, 81
\item Distances of points, 1
\item Lines of curvature of ellipsoid, 7
\item Pascal's theorem, 9
\end{itemize}

\section*{Analysis}
\begin{itemize}[leftmargin=*,noitemsep]
\item Determinants, 1, 9, 12, 52, 61, 69
\item Attractions and multiple integrals, 2, 3, 28, 29, 41, 44, 56, 63, 64, 88
\item Attraction of Ellipsoids, Legendre, 75; Jacobi, 89
\item Linear transformations, \&c.\ 12, 13, 14, 15, 16, 33, 34, 35, 39, 53, 54, 71, 74, 77, 100
\item Elliptic functions, 17, 18, 19, 21, 23, 24, 25, 33, 45, 57, 58, 62, 67, 87, 88, 93, 94, 99
\end{itemize}

\cleardoublepage

% ============================================
% MAIN MATTER: Pages 23--52 (Papers 1--6, start of 7)
% ============================================
\mainmatter

% --------------------------------------------
% PAPER 1: Pages 23--26
% --------------------------------------------
\chapter{On a Theorem in the Geometry of Position}
\label{chap:paper1}

\begin{center}
\small[From the Cambridge Mathematical Journal, vol.~II.\ (1841), pp.~267--271.]
\end{center}

\section*{}
\markboth{On a Theorem in the Geometry of Position}{}

We propose to apply the following (new?) theorem to the solution of two problems in Analytical Geometry.

Let the symbols
\[
\begin{array}{c}
\alpha, \; \beta \\
\alpha', \; \beta', \; \gamma' \\
\alpha'', \; \beta'', \; \gamma'', \; \delta'' \\
\&\mathrm{c.}
\end{array}
\]
denote the quantities
\begin{multline*}
\alpha, \quad \alpha\beta' - \alpha'\beta, \quad \alpha\beta'\gamma'' - \alpha\beta''\gamma' + \alpha'\beta''\gamma - \alpha'\beta\gamma'' \\
+ \alpha''\beta\gamma' - \alpha''\beta'\gamma, \quad \&\mathrm{c.}
\end{multline*}
(the law of whose formation is tolerably well known, but may be shown in a more convenient form than is commonly given to it).

Then, if
\[
\begin{array}{ccccc}
x_1, & y_1, & z_1, & w_1, & \&\mathrm{c.} \\
x_2, & y_2, & z_2, & w_2, & \&\mathrm{c.} \\
x_3, & y_3, & z_3, & w_3, & \&\mathrm{c.} \\
\&\mathrm{c.} & \&\mathrm{c.} & \&\mathrm{c.}
\end{array}
\]
denote any quantities whatever, we have the theorem
\[
\left|\begin{array}{cccc}
x_1 & y_1 & z_1 & \cdots \\
x_2 & y_2 & z_2 & \cdots \\
x_3 & y_3 & z_3 & \cdots \\
\vdots & \vdots & \vdots & \ddots
\end{array}\right|
\times
\left|\begin{array}{cccc}
\alpha_1 & \beta_1 & \gamma_1 & \cdots \\
\alpha_2 & \beta_2 & \gamma_2 & \cdots \\
\alpha_3 & \beta_3 & \gamma_3 & \cdots \\
\vdots & \vdots & \vdots & \ddots
\end{array}\right|
= \Sigma \left\{
\left|\begin{array}{cccc}
x_1\alpha_{\theta} + y_1\beta_{\theta} + \cdots &
x_1\alpha_{\phi} + y_1\beta_{\phi} + \cdots & \cdots \\
x_2\alpha_{\theta} + y_2\beta_{\theta} + \cdots &
x_2\alpha_{\phi} + y_2\beta_{\phi} + \cdots & \cdots \\
\vdots & \vdots & \ddots
\end{array}\right|
\right\},
\]
where $\theta, \phi, \ldots$ denote any permutation of $1, 2, 3, \ldots$.

\textit{(This theorem admits of a generalisation which we shall not have occasion to make use of, and which therefore we may notice at another opportunity.)}

\subsection*{Problem I}

To find the relation that exists between the distances of five points in space.

We have, in general, whatever $x_1, y_1, z_1, w_1$, \&c.\ denote,
\[
\begin{vmatrix}
x_1^2 + y_1^2 + z_1^2 + w_1^2 & -2x_1 & -2y_1 & -2z_1 & -2w_1 & 1 \\
x_2^2 + y_2^2 + z_2^2 + w_2^2 & -2x_2 & -2y_2 & -2z_2 & -2w_2 & 1 \\
x_3^2 + y_3^2 + z_3^2 + w_3^2 & -2x_3 & -2y_3 & -2z_3 & -2w_3 & 1 \\
\vdots & \vdots & \vdots & \vdots & \vdots & \vdots
\end{vmatrix} = 0.
\]

\noindent
Write $12$ for the distance between points 1 and 2, \&c.; then expanding, we obtain
\[
\begin{vmatrix}
0 & 12^2 & 13^2 & 14^2 & 15^2 & 1 \\
21^2 & 0 & 23^2 & 24^2 & 25^2 & 1 \\
31^2 & 32^2 & 0 & 34^2 & 35^2 & 1 \\
41^2 & 42^2 & 43^2 & 0 & 45^2 & 1 \\
51^2 & 52^2 & 53^2 & 54^2 & 0 & 1 \\
1 & 1 & 1 & 1 & 1 & 0
\end{vmatrix} = 0,
\]
which is easily expanded, though from the mere number of terms the process is somewhat long.

Precisely the same investigation is applicable to the case of four points in a plane, or three points in a straight line. Thus the former gives
\[
\begin{vmatrix}
0 & 12^2 & 13^2 & 14^2 & 1 \\
21^2 & 0 & 23^2 & 24^2 & 1 \\
31^2 & 32^2 & 0 & 34^2 & 1 \\
41^2 & 42^2 & 43^2 & 0 & 1 \\
1 & 1 & 1 & 1 & 0
\end{vmatrix} = 0.
\]

The latter gives
\[
\begin{vmatrix}
0 & 12^2 & 13^2 & 1 \\
21^2 & 0 & 23^2 & 1 \\
31^2 & 32^2 & 0 & 1 \\
1 & 1 & 1 & 0
\end{vmatrix} = 0.
\]

\subsection*{Problem II}

To find the relation that exists between the distances of five points on a sphere or four points in a circle.

The (additional) equation that exists between the distances of five points on a sphere or four points in a circle, has such a remarkable analogy with the preceding, that they almost require to be noticed at the same time.

\cleardoublepage

% --------------------------------------------
% PAPER 2: Pages 27--34
% --------------------------------------------
\chapter{On the Properties of a Certain Symbolical Expression}
\label{chap:paper2}

\begin{center}
\small[From the Cambridge Mathematical Journal, vol.~iii.\ (1841), pp.~62--71.]
\end{center}

\section*{}
\markboth{On the Properties of a Certain Symbolical Expression}{}

The series
\[
S = \frac{1}{2^s \cdot 1 \cdot 2 \cdots s} \left(\frac{d}{da\, da'} + \frac{d}{db\, db'} + \&\mathrm{c.}\right)^s
(aa' + bb' + \&\mathrm{c.} - n\,\mathrm{terms})^{i+s},
\]
possesses some remarkable properties, which it is the object of the present paper to investigate. We shall prove that the symbolical expression
\[
\frac{1}{2^s \cdot 1 \cdot 2 \cdots s} \left(\frac{d}{da\, da'} + \frac{d}{db\, db'} + \&\mathrm{c.}\right)^s
\]
is equivalent to
\[
\frac{(-1)^s}{2 \cdot 4 \cdots 2s} \left(\frac{d^2}{da\, da'} + \frac{d^2}{db\, db'} + \&\mathrm{c.}\right)^s,
\]
and that, putting
\[
\Delta = \frac{d^2}{da\, da'} + \frac{d^2}{db\, db'} + \&\mathrm{c.},
\]
the operation $S$ upon any homogeneous function $U$ of the order $2s$ gives
\[
S \cdot U = \frac{(-1)^s \cdot 2i \cdot (2i+2) \cdots (2i+2s-2)}{2^s \cdot 1 \cdot 2 \cdots s} U.
\]

Considering the expression
\[
\frac{1}{2^s \cdot 1 \cdot 2 \cdots s} \left(\frac{d}{da\, da'} + \&\mathrm{c.}\right)^s (aa' + bb' + \&\mathrm{c.})^{i+s},
\]
if for a moment we write
\[
(l+1)a' = a_1', \quad \&\mathrm{c.}; \quad \lambda_1 = \frac{a_1'}{a'}, \quad \&\mathrm{c.};
\quad P_1 = a_1'^2 + b_1'^2 + \&\mathrm{c.},
\]
this becomes
\[
\frac{1}{2^s \cdot 1 \cdot 2 \cdots s} \left(\frac{d}{da\, d\lambda_1} + \&\mathrm{c.}\right)^s P_1^{i+s}.
\]

Now it is immediately seen that $\Delta_1 - \frac{d^2}{d\lambda_1^2} = \frac{2i'(2i'+2-n)}{P_1}$; from which we may deduce
\[
\Delta_1^q \cdot P_1^{i'} = \frac{2i' \cdot (2i'+2) \cdots (2i'+2q-2) \cdot (2i'+2-n) \cdots (2i'+2q-n)}{P_1^q}.
\]

Now $U$ representing any homogeneous function of the order $2s$, it is easily seen that
\[
\Delta \cdot P^{i'} U = 2i' \cdot (2i'+2s-n) \cdot P^{i'-1} U + P^{i'} \cdot \Delta U,
\]
and repeating continually the operation $\Delta$, observing that $\Delta U$, $\Delta^2 U$, \&c.\ are of the orders $2(s-1)$, $2(s-2)$, \&c., we at length arrive at
\begin{multline*}
\Delta^q \cdot P^{i'} U = \frac{2i' \cdot (2i'+2q-4s-n)}{P_1} \Delta^{q-1} U \\
+ \frac{2i' \cdot (2i'+2) \cdot (2i'+2q-4s-n) \cdot (2i'+2q-4s-n-2)}{P_1^2} \Delta^{q-2} U \\
+ \cdots \cdots \\
+ 2i' \cdot (2i'+2) \cdots (2i'+2q) \cdot (2i'+2q-4s-n) \cdots (2i'-2s-n+2) \cdot P^{i'+q-s} U.
\end{multline*}

Substituting for the quantities involved in this expression, and putting, for simplicity,
\[
2i + 2 - n = 2i',
\]
we have, without any further reduction, except that
\[
\Delta^s \cdot P^{i+s} U = \frac{(-1)^s \cdot 2i \cdot (2i+2) \cdots (2i+2s-2)}{2^s \cdot 1 \cdot 2 \cdots s} U,
\]
or we have simply
\[
S \cdot U = \frac{(-1)^s \cdot 2i \cdot (2i+2) \cdots (2i+2s-2)}{2^s \cdot 1 \cdot 2 \cdots s} U,
\]
where
\[
\frac{d}{dt} = \frac{d}{da\, da'} + \frac{d}{db\, db'} + \&\mathrm{c.},
\]
so that $S = f\left(\frac{d}{dt}\right)$.

\textit{[2.]} Being determined by the equation
\[
\frac{1}{2^s \cdot 1 \cdot 2 \cdots s} \left(\frac{d}{da\, da'} + \&\mathrm{c.}\right)^s (aa' + bb' + \&\mathrm{c.})^{i+s}
= \frac{(-1)^s}{2 \cdot 4 \cdots 2s} \left(\frac{d^2}{da\, da'} + \&\mathrm{c.}\right)^s,
\]
or, as it may otherwise be written,
\[
\{aa' + bb' + \&\mathrm{c.}\}^{i+s}
= \frac{(-1)^s}{2 \cdot 4 \cdots 2s} \left(\frac{d^2}{da\, da'} + \&\mathrm{c.}\right)^s (aa' + bb' + \&\mathrm{c.})^{i+s}.
\]

\textit{[3.]} Or the function $\phi(a-x,\; b-y,\ldots)$ not becoming infinite within the limits of the integration, we have
\[
\iint \cdots \phi(a-x,\; b-y,\ldots)\, dx\, dy \ldots
= \phi\left(\frac{d}{da},\; \frac{d}{db},\ldots\right) \iint \cdots e^{-(ax + by + \ldots)} dx\, dy \ldots
\]
since, as applied to the function $\phi$, $\frac{d^2}{da^2} + \frac{d^2}{db^2} + \&\mathrm{c.}$ is equivalent to $0$; we have therefore
\[
\iint \cdots \phi(a-x,\; b-y,\ldots)\, dx\, dy \ldots
= \frac{\pi^{\frac{n}{2}}}{\Gamma\left(\frac{n}{2}\right)} \int \phi\left(\sqrt{a^2 + b^2 + \&\mathrm{c.}}\right) r^{n-1}\, dr.
\]

\cleardoublepage

% --------------------------------------------
% PAPER 3: Pages 35--40
% --------------------------------------------
\chapter{On Certain Definite Integrals}
\label{chap:paper3}

\begin{center}
\small[From the Cambridge Mathematical Journal, vol.~iii.\ (1841), pp.~138--144.]
\end{center}

\section*{}
\markboth{On Certain Definite Integrals}{}

In the first place, we shall consider the integral
\[
V = \iint \cdots e^{-(ax^2 + 2hxy + by^2 + \ldots)}\, dx\, dy \ldots,
\]
taken between the limits $-\infty$ and $+\infty$ for each of the variables $x, y, \ldots$.

\textit{[1.]} Let the variable $x$, on the second side of the equation, be replaced by $\xi$, where we have without difficulty
\[
\frac{dV}{da} = \frac{h^2}{a^2} \cdot \frac{\pi^{\frac{1}{2}}}{a^{\frac{1}{2}}} \cdot \frac{dV'}{db'},
\]
where $V'$ denotes the same integral as $V$, but with the coefficients modified as above; and we hence obtain
\[
\frac{dV}{da} = \frac{h^2}{a^2} V'.
\]

\textit{[2.]} Hence the function $f_0\, e^{-ax^2}\, dx$ (observing that $m = n = 0$) is given by the equation
\[
V = \frac{\pi^{\frac{1}{2}}}{a^{\frac{1}{2}}}, \quad \text{or} \quad V = \frac{\pi^{\frac{1}{2}}}{\sqrt{a}}.
\]

In general if $u$ be any function of $\xi$,
\[
\frac{d^r u}{da^r} = \frac{1}{2^r} \cdot \frac{d^r u}{d\xi^r},
\]
from which the values of the second side for any particular value of $u$ may be at once obtained.

\textit{[3.]} Substituting this value, also
\[
\frac{h^2}{a^2} = h'^2, \quad \frac{b}{a} = b', \quad \&\mathrm{c.},
\]
and $f(\xi^m)$ in the value of $V$, and observing that
\[
m = \frac{1}{2}, \quad n = \frac{3}{2}, \quad \&\mathrm{c.},
\]
we obtain the value of the integral in a series, which may be thus written:
\[
V = \frac{\pi^{\frac{n}{2}}}{\sqrt{\Delta}} \left\{1 + \frac{1}{2} \cdot \frac{d}{d\xi} + \frac{1 \cdot 3}{2 \cdot 4} \cdot \frac{d^2}{d\xi^2} + \&\mathrm{c.}\right\} f(\xi^m),
\]
where $\Delta$ denotes the determinant formed from the coefficients $a, h, b, \ldots$.

\cleardoublepage

% --------------------------------------------
% PAPER 4: Pages 41--46
% --------------------------------------------
\chapter{On Certain Expansions, in series of Multiple Sines and Cosines}
\label{chap:paper4}

\begin{center}
\small[From the Cambridge Mathematical Journal, vol.~iii.\ (1842), pp.~162--167.]
\end{center}

\section*{}
\markboth{On Certain Expansions, in series of Multiple Sines and Cosines}{}

In the following paper I propose to consider certain expansions in series of multiple sines and cosines, which present themselves in the investigation of various problems in pure and applied mathematics. The method which I employ is essentially the same as that made use of by Fourier, Poisson, and others; but the results are obtained in a more general form, and some remarkable properties of the coefficients are pointed out.

\textit{[1.]} Let $f(\theta)$ be a function which, between the limits $0$ and $\pi$, is capable of being expanded in a series of the form
\[
f(\theta) = a_0 + a_1 \cos \theta + a_2 \cos 2\theta + \cdots + a_r \cos r\theta + \cdots
\]
Then multiplying by $\cos r\theta$ and integrating from $0$ to $\pi$, we obtain
\[
a_r = \frac{2}{\pi} \int_0^{\pi} f(\theta) \cos r\theta\, d\theta,
\]
(for $r > 0$), and
\[
a_0 = \frac{1}{\pi} \int_0^{\pi} f(\theta)\, d\theta.
\]

\textit{[2.]} Now, between the limits $0$ and $\pi$, the function $f(e^{i\theta} - e^{-i\theta}) \cos r\theta \cdot (e^{i\theta} - e^{-i\theta})$ will always be expansible in a series of multiple cosines of $r\theta$; and the same is true of any rational function of $e^{i\theta}$, $e^{-i\theta}$. Hence, if we write
\[
\theta - \theta' = \Sigma - \sigma',
\]
and expand $f(e^{i(\theta-\theta')})$ in a series of multiple cosines of $r(\theta - \theta')$, we obtain
\[
\cos k(\theta - \theta') = \cos k(\Sigma - \sigma') \left\{\cos k(\theta - \sigma) \cos k(\theta' - \sigma') + \sin k(\theta - \sigma) \sin k(\theta' - \sigma')\right\}.
\]

By choosing for $f(e^{i\theta} - e^{-i\theta},\; e^{i\phi} - e^{-i\phi},\ldots)$ functions expansible without sines, or without cosines, various remarkable results may be obtained.

\cleardoublepage

% --------------------------------------------
% PAPER 5: Pages 47--49
% --------------------------------------------
\chapter{On the Intersection of Curves}
\label{chap:paper5}

\begin{center}
\small[From the Cambridge Mathematical Journal, vol.~iii.\ (1843), pp.~211--213.]
\end{center}

\section*{}
\markboth{On the Intersection of Curves}{}

The following theorem is quoted in a note of Chasles' Aper\c{c}u Historique, and was suggested to me by a theorem of Steiner's, which will be given hereafter.

Consider generally two curves, $U_m = 0$, $V_n = 0$, of the orders $m$ and $n$ respectively, and a curve of the $r^{\text{th}}$ order ($r$ not less than $m$ or $n$). The theorem asserts that through the $mn$ points of intersection of the two curves, there may be drawn
\[
\frac{r(r+3)}{2} - \frac{(r-m)(r-m+1)}{2} - \frac{(r-n)(r-n+1)}{2}
\]
curves of the $r^{\text{th}}$ order; and that the number of conditions to be satisfied by the curve of the $r^{\text{th}}$ order is
\[
mn - \frac{1}{2}(m + n - r - 1)(m + n - r - 2).
\]

For the number of terms in the general equation of the $r^{\text{th}}$ order is $\frac{1}{2}(r+1)(r+2)$; the curve passes through $\frac{1}{2}r(r+3)$ given points, and through the $mn - \frac{1}{2}(m+n-r-1)(m+n-r-2)$ remaining points of the intersection of the two curves $U_m = 0$, $V_n = 0$.

Such a curve passes through $\frac{1}{2}r(r+3)$ given points, and though the $mn - \frac{1}{2}(m+n-r)(m+n-r-1)$ conditions are not explicitly satisfied, the theorem shows that the curve does in fact pass through the remaining points of intersection.

\cleardoublepage

% --------------------------------------------
% PAPER 6: Pages 50--52 (begins; continues beyond page 52)
% --------------------------------------------
\chapter{On the Motion of Rotation of a Solid Body}
\label{chap:paper6}

\begin{center}
\small[From the Cambridge Mathematical Journal, vol.~iii.\ (1843), pp.~224--232.]
\end{center}

\section*{}
\markboth{On the Motion of Rotation of a Solid Body}{}

In the fifth volume of Liouville's Journal, there is a very elegant paper by M.~Ostrogradsky, on the motion of rotation of a solid body about a fixed point. The author deduces the equations of motion by means of the principle of virtual velocities, and obtains some remarkable integrals of these equations. The object of the present paper is to investigate the same problem by a different method, and to obtain some additional results.

Let $Ax$, $Ay$, $Az$ be the principal axes of the body at the fixed point; and let $A$, $B$, $C$ be the principal moments of inertia. Let $\lambda$, $\mu$, $\nu$ be the direction-cosines of a fixed line in space, with respect to the moving axes; and let $p$, $q$, $r$ be the angular velocities about these axes. Then the equations of motion are
\begin{align}
A \frac{dp}{dt} + (C - B)qr &= M g (\mu \cdot z_0 - \nu \cdot y_0), \\
B \frac{dq}{dt} + (A - C)rp &= M g (\nu \cdot x_0 - \lambda \cdot z_0), \\
C \frac{dr}{dt} + (B - A)pq &= M g (\lambda \cdot y_0 - \mu \cdot x_0),
\end{align}
where $x_0$, $y_0$, $z_0$ are the coordinates of the centre of gravity with respect to the principal axes.

Also, since $\lambda$, $\mu$, $\nu$ are the direction-cosines of a fixed line,
\begin{align}
\frac{d\lambda}{dt} &= \mu r - \nu q, \\
\frac{d\mu}{dt} &= \nu p - \lambda r, \\
\frac{d\nu}{dt} &= \lambda q - \mu p.
\end{align}

Suppose that $Ax$, $Ay$, $Az$ are referred to axes $A\xi$, $A\eta$, $A\zeta$, by the quantities $l$, $m$, $n$, $k$, analogous to $\lambda$, $\mu$, $\nu$; then we have
\[
\begin{aligned}
a &= \cos^2 l + \sin^2 l \cos \theta, \\
a' &= \cos l \cos l' + \sin l \sin l' \cos(\Sigma - \sigma'), \\
a'' &= \cos l \cos l'' + \sin l \sin l'' \cos(\Sigma + \sigma'),
\end{aligned}
\qquad
\begin{aligned}
\beta &= \cos^2 m + \sin^2 m \cos \theta, \\
\beta' &= \cos m \cos m' + \sin m \sin m' \cos(\mathrm{M} - \mu'), \\
\beta'' &= \cos m \cos m'' + \sin m \sin m'' \cos(\mathrm{M} + \mu'),
\end{aligned}
\]
and similar expressions for $\gamma$, $\gamma'$, $\gamma''$.

\vfill
\begin{center}
\textit{[Paper 6 continues beyond page 52]}
\end{center}

% ============================================
% END OF PAGES 1--52
% ============================================
\cleardoublepage

\end{document}
